\bookchapter{บทที่ 43: วิศวกรรมบทสนทนา—เทคนิคพูดคุยกับ AI อย่างมีประสิทธิภาพ}

\chapterintro{บทสนทนาคือรูปแบบการโต้ตอบหลักของการเขียนโปรแกรมด้วย AI แต่โปรดสังเกต: “บทสนทนา” ไม่เท่ากับ “คุยเล่น”—บทสนทนาเพื่อเขียนโปรแกรมกับ AI ที่มีประสิทธิภาพเป็นเทคนิคการสื่อสารที่มีโครงสร้าง เราเรียกมันว่า “วิศวกรรมบทสนทนา” มันศึกษาว่า: จะใช้รอบการสนทนาให้น้อยที่สุดอย่างไรเพื่อให้ได้โค้ดคุณภาพสูงสุด พูดตรง ๆ บทนี้คือหลักสูตรสอนวิธี “สื่อสารให้ชัด” กับ AI นั่นเอง}

\booksection{กลยุทธ์แบ่งชั้นของ Prompt}

บทสนทนากับ AI แบ่งได้สามชั้น แต่ละชั้นทำหน้าที่ต่างกัน พอเข้าใจสามชั้นนี้ คุณจะต่อ Prompt ที่มีประสิทธิภาพได้เหมือนเล่นบล็อกเสริม

\begin{center}
\begin{tabularx}{\linewidth}{|X|X|X|X|}
\hline
\textbf{ชั้น} & \textbf{หน้าที่} & \textbf{ตัวอย่าง} & \textbf{ความถี่} \tabularnewline
\hline
ชั้นระบบ & กำหนดบทบาทของ AI และข้อกำหนดรวมทั้งหมด & "คุณเป็นผู้เชี่ยวชาญ React โค้ดต้องเขียนด้วย TypeScript" & ตั้งไว้ครั้งเดียวตอนเริ่มโปรเจกต์ \tabularnewline
ชั้นงาน & อธิบายงานเฉพาะหน้าที่จะทำตอนนี้ & "สร้างคอมโพเนนต์ฟอร์ม มีช่องชื่อผู้ใช้และรหัสผ่าน" & ระบุตอนเริ่มแต่ละงาน \tabularnewline
ชั้นฟีดแบ็ก & ประเมินและแก้ไขผลงานของ AI & "ช่องรหัสผ่านยังขาดสวิตช์แสดง/ซ่อน ช่วยเพิ่มด้วย" & ให้ตามคุณภาพของผลงาน \tabularnewline
\hline
\end{tabularx}
\end{center}

\begin{center}\small\color{gray}สามชั้นของ Prompt ในบทสนทนา\end{center}

Prompt ชั้นระบบมักเขียนไว้ในไฟล์ความจำของโปรเจกต์ (เช่น CLAUDE.md) ไม่ต้องพูดซ้ำทุกรอบคุย ชั้นงานคือสิ่งที่คุณใช้บ่อยที่สุดในบทสนทนา—ควรมีสามองค์ประกอบ “ทำอะไร” “ใช้เทคโนโลยีอะไร” “มีข้อจำกัดอะไร” ส่วนชั้นฟีดแบ็กคือชั้นที่วัดฝีมือจริง—ฟีดแบ็กดีทำให้ AI แก้ได้แม่นยำ ฟีดแบ็กแย่ทำให้มันแค่เดาใหม่ไปวัน ๆ

\begin{tcbinfo}{เคล็ดลับ: สามจังหวะของฟีดแบ็ก}
เทมเพลตฟีดแบ็กที่ใช้ได้จริงคือ “ชี้ปัญหา + บอกสิ่งที่คาดหวัง + ให้บริบท” เช่น “การจัดการข้อผิดพลาดบรรทัด 25 ใช้แค่ console.log (ชี้ปัญหา) ในสภาพแวดล้อมจริงควรใช้ตัวจัดการข้อผิดพลาดกลางและคืนข้อความที่ผู้ใช้เข้าใจง่าย (บอกสิ่งที่คาดหวัง) ในโปรเจกต์มีฟังก์ชัน handleError ให้ใช้ (ให้บริบท)” ฟีดแบ็กแบบนี้ได้ผลกว่า “แก้การจัดการข้อผิดพลาดหน่อย” ถึงสิบเท่า

\end{tcbinfo}

\booksection{บทสนทนาดีบั๊ก: วิธีบรรยายข้อผิดพลาดให้ AI หาต้นตอได้ไว}

การดีบั๊กเป็นสถานการณ์ที่พบบ่อยที่สุดอย่างหนึ่งในการเขียนโปรแกรมด้วย AI เมื่อเจอบั๊ก วิธีที่คุณบรรยายปัญหาให้ AI ฟังจะกำหนดโดยตรงว่ามันช่วยคุณหาต้นตอและแก้ได้เร็วแค่ไหน นี่คือสี่ประเด็นสำคัญของบทสนทนาดีบั๊ก:

\begin{enumerate}
  \item ให้ AI "ดู" ข้อมูลต้นฉบับ ไม่ใช่คำบรรยายของคุณ ดึง error log, stack trace และโค้ดที่เกี่ยวข้องมาอ้างอิงด้วย @ ให้ AI ตรง ๆ อย่าบรรยายแบบ "เหมือนมีอะไร error นะ" AI เห็นข้อมูลต้นฉบับแล้วแม่นยำกว่าเห็นความเข้าใจของคุณ
  \item อธิบายช่องว่างระหว่าง "พฤติกรรมที่คาดหวัง" กับ "พฤติกรรมที่เกิดจริง" "ฉันคาดว่ากดปุ่มแล้วจะไปหน้าแรก (คาดหวัง) แต่จริง ๆ กดแล้วหน้าไม่มีปฏิกิริยาอะไรเลย (เกิดจริง)"—คำบรรยายแบบนี้มีประโยชน์กว่า "ปุ่มพัง" มาก
  \item ให้ขั้นตอนทำซ้ำ "พิมพ์ชื่อผู้ใช้ก่อน แล้วกดปุ่มเข้าสู่ระบบ จะเจอข้อผิดพลาดนี้"—ขั้นตอนทำซ้ำช่วยให้ AI จำลองเส้นทางการใช้งานของคุณได้ หาต้นตอปัญหาง่ายขึ้น
  \item บอกให้ชัดว่าคุณลองอะไรมาแล้ว "ฉันเช็ก network request แล้ว คืน 200 ข้อมูลปกติ"—สิ่งนี้ช่วยให้ AI ตัดทิศทางที่ตรวจแล้วออกไป ไม่เปลืองเวลาไล่ซ้ำ
\end{enumerate}

คำบรรยายปัญหาที่แย่

\begin{lstlisting}
กดปุ่มแล้วไม่มีอะไรเกิดขึ้น ช่วยดูหน่อย
\end{lstlisting}

คำบรรยายปัญหาที่ดี

\begin{lstlisting}
กดปุ่มเข้าสู่ระบบแล้วหน้าไม่ตอบสนอง พฤติกรรมที่คาดหวัง: กดแล้วไปที่หน้า /home พฤติกรรมที่เกิดจริง: กดแล้วไม่มีปฏิกิริยาใด ๆ คอนโซลไม่มี error ขั้นตอนทำซ้ำ: 1. ในหน้าเข้าสู่ระบบ พิมพ์ชื่อผู้ใช้ admin รหัสผ่าน 123456 2. กดปุ่ม "เข้าสู่ระบบ" 3. หน้าไม่ตอบสนอง สิ่งที่ตรวจแล้ว: - network request ปกติ คืนค่า 200 - เหตุการณ์ onClick ผูกไว้แล้ว - เช็กโค้ดคอมโพเนนต์ LoginButton แล้ว ไฟล์ที่เกี่ยวข้อง: @src/components/LoginButton.tsx @src/pages/Login.tsx
\end{lstlisting}

\booksection{กลยุทธ์วนซ้ำ: ก้าวเล็กวิ่งไว vs ปรับใหญ่ทีเดียว}

การวนซ้ำในการเขียนโปรแกรมด้วย AI มีสองกลยุทธ์พื้นฐาน เหมาะกับสถานการณ์ต่างกัน เข้าใจความต่างจะช่วยให้คุณเลือกถูกในแต่ละสถานการณ์

กลยุทธ์ “ก้าวเล็กวิ่งไว” เหมาะกับการพัฒนาฟีเจอร์และแก้บั๊ก แก่นคือ: แต่ละครั้งให้ AI แก้เล็กน้อยแค่จุดเดียว ตรวจผ่านแล้วค่อยทำจุดถัดไป ข้อดีคือทุกก้าวตรวจได้ ย้อนกลับได้ ความเสี่ยงควบคุมได้; ข้อเสียคือช้ากว่า ต้องโต้ตอบถี่ เหมาะกับ: พัฒนาฟีเจอร์ใหม่ สำรวจเมื่อยังไม่แน่ใจทางที่ดีที่สุด และวนซ้ำอย่างระมัดระวังเมื่อยังไม่ค่อยวางใจผลงานของ AI

กลยุทธ์ “ปรับใหญ่ทีเดียว” เหมาะกับการแก้ไขขนาดใหญ่ที่เป้าหมายชัดเจน แก่นคือ: อธิบายเป้าหมายฉบับเต็มครั้งเดียว ให้ AI วางแผนและลงมือทำการแก้ทั้งหมดด้วยตัวเอง ข้อดีคือเร็ว AI เห็นภาพรวมจึงตัดสินใจเชิงสถาปัตยกรรมได้ดีกว่า; ข้อเสียคือเสี่ยงสูง ถ้า AI เข้าใจเพี้ยน ค่าแก้ไขจะแพง เหมาะกับ: อัปเกรดสถาปัตยกรรมที่ชัดเจน ย้ายโค้ดขนาดใหญ่ และ refactor ที่มีชุดทดสอบรองรับพอ

\begin{center}
\begin{tabularx}{\linewidth}{|X|X|X|}
\hline
\textbf{มิติ} & \textbf{ก้าวเล็กวิ่งไว} & \textbf{ปรับใหญ่ทีเดียว} \tabularnewline
\hline
ปริมาณการแก้ต่อครั้ง & หนึ่งฟีเจอร์หรือหนึ่งบั๊ก & หนึ่งโมดูลหรือหนึ่งชั้นสถาปัตยกรรม \tabularnewline
ความถี่การตรวจ & ตรวจทันทีหลังแก้ทุกครั้ง & แก้เสร็จหมดแล้วตรวจรวมครั้งเดียว \tabularnewline
จำนวนรอบสนทนา & มาก (5-10 รอบ/ฟีเจอร์) & น้อย (1-3 รอบ/โมดูล) \tabularnewline
ความเสี่ยง & ต่ำ ย้อนกลับได้ทีละก้าว & สูง ค่าแก้ความเพี้ยนแพง \tabularnewline
สถานการณ์ที่เหมาะ & ฟีเจอร์ใหม่ งานสำรวจ & refactor หรือย้ายที่เป้าหมายชัด \tabularnewline
เงื่อนไขเบื้องต้น & ไม่มีข้อกำหนดพิเศษ & ต้องมีชุดทดสอบครอบคลุม \tabularnewline
\hline
\end{tabularx}
\end{center}

\begin{center}\small\color{gray}เปรียบเทียบกลยุทธ์วนซ้ำสองแบบ\end{center}

ในการทำจริง สองกลยุทธ์มักใช้คู่กัน เช่น ตอน refactor หนึ่งโมดูล เริ่มด้วย “ปรับใหญ่ทีเดียว” ให้ AI เสนอแผนรวม ตรวจแผนแล้วค่อยใช้ “ก้าวเล็กวิ่งไว” ลงมือทีละขั้น วิธีนี้ได้มุมมองรวมและคุมความเสี่ยงได้พร้อมกัน

\booksection{การจัดการบริบทในบทสนทนาหลายรอบ}

ยิ่งคุยหลายรอบ คุณยิ่งต้องจัดการบริบทการสนทนาอย่างตั้งใจ เพื่อเลี่ยงปัญหา “AI ลืมข้อมูลช่วงแรก ๆ” นี่คือเทคนิคจัดการบริบทที่ใช้ได้จริง

หนึ่ง: “ตอกหมุด” ข้อมูลสำคัญ ตอนเริ่มบทสนทนาเขียนข้อจำกัดสำคัญให้ชัด และย้ำเป็นระยะในรอบถัด ๆ ไป เช่น “จำไว้นะ เราเก็บข้อมูลใน localStorage ใช้ฐานข้อมูลไม่ได้”—ย้ำข้อจำกัดนี้ในคำขอแก้ไขที่เกี่ยวข้อง กัน AI หลงลืม

สอง: ใช้เทคนิค “สรุปแล้วต่อ” เมื่อบทสนทนายาวมาก ขอให้ AI “สรุปความคืบหน้าและสิ่งที่ค้างอยู่” แล้วเปิดบทสนทนาใหม่โดยใช้สรุปนั้นเป็นจุดตั้งต้น วิธีนี้เก็บบริบทสำคัญไว้และเลี่ยงการเสื่อมของบทสนทนายาว

สาม: “หนึ่งบทสนทนา หนึ่งงาน” นี่คือกลยุทธ์จัดการบริบทที่ง่ายที่สุดและได้ผลที่สุด แต่ละบทสนทนารับงานเดียวหรือฟีเจอร์เดียว เสร็จแล้วเปิดบทสนทนาใหม่ทำงานถัดไป มันจำกัดความยาวบทสนทนาโดยธรรมชาติ บริบทจึงคงชัดเจน

สี่: ใช้ฟีเจอร์จัดการบริบทของเครื่องมือ เครื่องมือเขียนโปรแกรม AI ส่วนใหญ่มีปุ่ม “ล้างบริบท” หรือ “เซสชันใหม่”—เมื่อรู้สึกว่าคุณภาพคำตอบของ AI ทรุด อย่าลังเล เปิดเซสชันใหม่เลย เสียประวัติบทสนทนาไปบ้างยังดีกว่าฝืนต่อในบริบทที่เสื่อมแล้ว

\begin{tcbpractice}{ลงมือทำ: สวมบท "คนไข้" ให้ AI สักครั้ง}
จงใจสร้างหน้าเล็ก ๆ ในเบราว์เซอร์ แล้ววางบั๊กไว้หนึ่งตัว (เช่น กดปุ่มแล้วไม่มีปฏิกิริยา) จากนั้นอธิบายปัญหาให้ AI ครบชุดตามรูปแบบ “คำบรรยายปัญหาที่ดี” ในบทนี้: ข้อมูลต้นฉบับ, คาดหวัง vs เกิดจริง, ขั้นตอนทำซ้ำ, สิ่งที่ตรวจแล้ว ลองสัมผัสความต่างของความเร็วที่ AI จับปัญหาเจอ—คุณจะเข้าใจว่าทำไม “การบรรยายปัญหา” ถึงเป็นศาสตร์อย่างหนึ่งเสียเอง

\end{tcbpractice}

\begin{tcbsummary}{สรุปบทนี้}
\begin{itemize}
  \item Prompt สามชั้น: ชั้นระบบ (เขียนในไฟล์ความจำ), ชั้นงาน (ทำอะไร+เทคโนโลยี+ข้อจำกัด), ชั้นฟีดแบ็ก (ชี้ปัญหา+คาดหวัง+บริบท)
  \item สี่ประเด็นดีบั๊ก: ให้ข้อมูลต้นฉบับ, คาดหวัง vs เกิดจริง, ขั้นตอนทำซ้ำ, สิ่งที่ตรวจแล้ว
  \item ก้าวเล็กวิ่งไว: เสี่ยงต่ำแต่รอบเยอะ เหมาะกับฟีเจอร์ใหม่; ปรับใหญ่ทีเดียว: เสี่ยงสูงแต่รอบน้อย ต้องมีชุดทดสอบรองรับ
  \item สี่เทคนิคจัดการบริบท: ตอกหมุดข้อมูลสำคัญ, สรุปแล้วต่อ, หนึ่งบทสนทนาหนึ่งงาน, เปิดเซสชันใหม่เมื่อจำเป็นโดยไม่ลังเล
\end{itemize}

\end{tcbsummary}
